\chapter{Mechanical Derivation of the FCC Lattice Constants}
\label{app:fcc_constants}
\addtocontents{toc}{\protect\setcounter{tocdepth}{0}}

\textbf{Subject:} The Mechanical Constitution of the Soft Floor and the Dual Nature of the Constants.

\section{The Soft Floor Node Density ($\rho_D$) — The Potential Reservoir}
With the Master Geometric Law ($\rho(r) = \mathcal{F}_i \cdot \rho_S \left(\frac{\ell_S}{r}\right)^4$) and the Spherical Form Factor ($\mathcal{F}_D = 1.5$) strictly defined in \hyperref[app:conservation_action]{Appendix B}, we can directly derive the volumetric stiffness density of the Soft Floor node by evaluating the law at the Soft Floor Coherence Scale ($r_D$).

\begin{itemize}
    \item \textbf{Definition:} $\rho_D$ represents the Operational Density of the vacuum node that stores the potential energy for quantum action.
    \item \textbf{Calculation:} Using $r_D \approx 7.8900 \times 10^{-27} \text{ m}$:
\end{itemize}
\begin{equation}
    \rho_D = 1.5 \cdot \rho_S \cdot \left(\frac{\ell_S}{r_D}\right)^4 \approx 1.3615 \times 10^{62} \text{ kg/m}^3
\end{equation}

\textit{(Note: This density is the mechanical source of the quantum action $h$, as mapped in the Fractal Architecture of Appendix E).}

\section{The Soft Floor Torsional Invariant ($J_D$) — The Master Coupling Gear}
As established in \hyperref[app:conservation_action]{Appendix B}, the Universal Torsional Invariant ($J_i = \rho_i r_i^4$) is the conserved mechanical quantity that governs topological phase transitions. Having derived both the scale ($r_D$) and the operational density ($\rho_D$) of the Soft Floor, we can define its specific invariant.

\begin{itemize}
    \item \textbf{Definition:} $J_D$ is the irreducible torsional resistance of a single Soft Floor spherical node. It is the central mechanical gear that projects the Substrate’s latent potential into the observable action ($h$) and gravitational coupling ($G$) of the active vacuum.
    \item \textbf{Calculation:}
\end{itemize}
\begin{equation}
    J_D = \rho_D \cdot r_D^4
\end{equation}
\begin{equation}
    J_D \approx (1.3615 \times 10^{62}) \cdot (7.8900 \times 10^{-27})^4 \approx 5.2763 \times 10^{-43} \text{ kg}\cdot\text{m}
\end{equation}

\textit{(Note: The strict mathematical proof demonstrating how $J_D$ mechanically generates the exact empirical values of $h$ and $G$ is detailed in Appendix E).}

\section{Vacuum Tension ($\tau_D$) — The Dual-Faced Stiffness}
The vacuum exhibits two distinct moduli depending on the mode of propagation: Shear (Transverse) and Bulk (Longitudinal).

\subsection{Surface Stiffness (Shear Modulus)}
The propagation of light is governed by the surface interface of the node. We define $\tau_D$ as the fundamental surface tension constant:
\begin{equation}
    \tau_D = \frac{c_E^4}{G \cdot r_D} \approx 1.5365 \times 10^{70} \text{ N/m}
\end{equation}
\begin{itemize}
    \item \textbf{Role:} When the vacuum transmits a transverse wave (Light), $\tau_D$ acts strictly as the shear modulus. This elastic snap-back of the lattice skin drives the propagation speed $c_E$.
\end{itemize}

\subsection{Volume Stiffness (Bulk Modulus)}
The background stress that supports the universe against structural failure and mediates static gravity is volumetric. We define the volumetric Gravitational Yield Tension ($\tau_{D,bulk}$) of the FCC lattice:
\begin{equation}
    \tau_{D,bulk} = \frac{\tau_D}{r_D} = \frac{c_E^4}{G \cdot r_D^2} \approx 1.9474 \times 10^{96} \text{ N/m}^2
\end{equation}
\begin{itemize}
    \item \textbf{Role:} Acts as the Incompressible Bulk Modulus. It is the ultimate structural yield strength of the Soft Floor lattice.
\end{itemize}

\section{Vacuum Inertia ($\mu_D$) and Mass Discretization}
Inertia is the resistance to change. We postulate that the vacuum possesses an intrinsic node interface inertia ($\mu_D$), defined by the contiguous 2D surface density of the Soft Floor nodes.
\begin{equation}
    \mu_D = \frac{c_E^2}{G \cdot r_D} \approx 1.7095 \times 10^{53} \text{ kg/m}^2
\end{equation}
This single constant governs both Light and Mass depending on the dimension of interaction.

\subsection{The Surface Interaction (Speed of Light)}
Light is a transverse shear wave surfing the 2D interface of the contiguous nodes. It encounters $\mu_D$ directly as an areal surface load:
\begin{equation}
    c_E = \sqrt{\frac{\tau_D}{\mu_D}}
\end{equation}
\textit{(Verification: $\sqrt{1.5365 \times 10^{70} / 1.7095 \times 10^{53}} \approx 2.9979 \times 10^8 \text{ m/s}$).}

\subsection{The Volumetric Projection (Holographic Drag)}
\label{subsec:volumetric_projection}

Matter is a volumetric displacement. A particle vibrating in the lattice displaces the interstitial condensate in 3D volume. Therefore, it encounters the lattice surface inertia ($\mu_D$) projected into 3D space.

To calculate macroscopic mass, we must mathematically map this 2D planar impedance into a 3D volumetric density. We define the effective interstitial vacuum condensate ($\rho_{IC}$) by projecting the surface inertia via the Dimensional Projection Constant ($\iota_V \equiv 1 \text{ m}^{-1}$):
\[
\rho_{IC} = \mu_D \cdot \iota_V \approx 1.7095 \times 10^{53} \text{ kg/m}^3
\]

\textbf{Mechanism:} Mass ($m$) is strictly the hydrodynamic drag of this projected condensate:
\[
m = C_d \cdot \rho_{IC} \cdot V_{\text{eff}} \cdot \sin(\Delta\phi)
\]

\textbf{Conclusion:} Light propagates transversely and feels the 2D skin ($\mu_D$). Mass displaces volumetrically and feels the 3D bulk ($\rho_{IC}$). They are the exact same physical entity, separated strictly by dimensional projection.

\subsection{The Quantum Mass ($M_D$) — The Node Gear}
While $\mu_D$ defines the 2D surface load of the vacuum interface, we must define the absolute localized 3D mass ($M_D$) of a single Soft Floor node.

\vspace{0.4cm}
\noindent \textbf{Derivation 1: The Geometric Origin (Volume)}\\
Because the Soft Floor node is a physical structure, its quantum mass is simply the Active Vacuum Density ($\rho_D$) multiplied by its fundamental 3D spherical volume:
\begin{equation}
    M_D = \rho_D \cdot \left(\frac{4}{3}\pi r_D^3\right) \approx 2.799 \times 10^{-16} \text{ kg}
\end{equation}

\vspace{0.4cm}
\noindent \textbf{Derivation 2: The Kinematic Origin (Quantum of Action)}\\
As established by the fractal geometry in Appendix E, the fundamental unreduced action ($h$) of the vacuum is generated by the geometric volumetric lever ($\frac{4}{3}\pi \cdot \rho_D \cdot r_D^4$) operating at the kinematic speed limit ($c_E$):
\begin{equation}
    h = \left[ \frac{4}{3}\pi \cdot \rho_D \cdot r_D^3 \right] \cdot c_E \cdot r_D = M_D \cdot c_E \cdot r_D
\end{equation}

This confirms that the physical spherical mass ($M_D$) is the exact inertial generator of the unreduced wave action ($h$):
\begin{equation}
    M_D = \frac{h}{c_E \cdot r_D}
\end{equation}

Because standard quantum mechanics defines localized particle states using the reduced Planck constant ($\hbar = h / 2\pi$), we can immediately translate this kinematic relationship to the localized mass. Substituting $h = 2\pi\hbar$ into the wave equation yields:
\begin{equation}
    M_D = \frac{2\pi \hbar}{c_E \cdot r_D}
\end{equation}

\textbf{Conclusion:} The unreduced wave action ($h$) and the localized particle action ($\hbar$) are kinematic duals that isolate the exact same physical spherical mass. The emergence of the $2\pi$ multiplier in the $\hbar$ equation strictly preserves the rotational topology of the localized node, ensuring the physical geometry of the vacuum remains entirely spherical.

\subsection{The Fundamental Drag Constant ($K_D$)}
$K_D$ is the linear mass gradient of the vacuum node. It defines the direct mathematical conversion between a particle's spatial vibration amplitude ($r_{amp}$) and its resulting hydrodynamic mass ($m$).

\textbf{Calculation:}
\begin{equation}
    K_D = \rho_{IC} \cdot (4\pi r_D^2)
\end{equation}
\begin{equation}
    K_D \approx (1.7095 \times 10^{53}) \cdot (7.8228 \times 10^{-52}) \approx 133.77 \text{ kg/m}
\end{equation}

\textbf{Role:} For every $1 \text{ meter}$ of volumetric expansion (amplitude) a particle forces onto the vacuum condensate, the vacuum exerts $133.77 \text{ kg}$ of inertial resistance ($m = K_D \cdot r_{amp}$).

\section{Summary of Soft Floor Constants}
The ``D'' constants define the constitution and absolute mechanical boundaries of the observable vacuum.

\begin{table}[htbp]
\centering
\caption{Summary of the Soft Floor Vacuum Constants}
\vspace{0.2cm}
\label{tab:soft_floor_constants}
\renewcommand{\arraystretch}{1.5}
\setlength{\tabcolsep}{4pt}
\small
\begin{tabular}{@{}llcp{0.45\textwidth}@{}}
\toprule
\textbf{Property} & \textbf{Symbol} & \textbf{Value (SI)} & \textbf{TFT Mechanical Definition} \\
\midrule
Operational Density & $\rho_D$ & $\approx 1.3615 \times 10^{62} \text{ kg/m}^3$ & Baseline mechanical mass-density of the active FCC node (source of $h$). \\
Torsional Invariant & $J_D$ & $\approx 5.2763 \times 10^{-43} \text{ kg}\cdot\text{m}$ & The absolute torsional resistance of the node (mediates $h$ and $G$). \\
Shear Tension & $\tau_D$ & $\approx 1.5365 \times 10^{70} \text{ N/m}$ & Lattice Surface Tension (Source of light speed). \\
Gravitational Yield Tension & $\tau_{D,bulk}$ & $\approx 1.9474 \times 10^{96} \text{ N/m}^2$ & Lattice Gravitational Yield Strength. \\
Node Interface Inertia & $\mu_D$ & $\approx 1.7095 \times 10^{53} \text{ kg/m}^2$ & Emergent Inertial Density for FCC lattice. \\
Interstitial Condensate & $\rho_{IC}$ & $\approx 1.7095 \times 10^{53} \text{ kg/m}^3$ & Interstitial Condensate Density (Holographic Drag). \\
Quantum Mass & $M_D$ & $\approx 2.799 \times 10^{-16} \text{ kg}$ & 3D localized spherical inertia of a single quantum node. \\
Fundamental Drag & $K_D$ & $\approx 133.77 \text{ kg/m}$ & Absolute Hydrodynamic Coupling Constant. \\
\bottomrule
\end{tabular}
\end{table}